\chapter*{ЗАКЛЮЧЕНИЕ}
\addcontentsline{toc}{chapter}{ЗАКЛЮЧЕНИЕ}

В рамках магистерской работы разработана и экспериментально проверена платформа видеоаналитики реального времени для мониторинга охраняемых зон и событийного анализа. Основной результат работы состоит в создании контура «обнаружение -- сопровождение -- событийный анализ -- надежная доставка», в котором алгоритмическая обработка видеопотока связана с транзакционной фиксацией событий, асинхронной публикацией сведений о событиях внешним потребителям и проверяемой ролевой моделью доступа.

По результатам исследования получены следующие основные результаты:
\begin{enumerate}
    \item Проанализирована предметная область охранного видеонаблюдения, рассмотрены существующие подходы к построению платформ видеоаналитики и сформулированы функциональные и нефункциональные требования к разрабатываемому решению.
    \item Обоснован выбор методов обнаружения, сопровождения и правиловой генерации событий: одноэтапный детектор семейства YOLO, сопровождение по результатам обнаружения на основе SORT/DeepSORT и набор интерпретируемых событийных правил \texttt{zone\_enter}, \texttt{line\_cross}, \texttt{dwell\_time}.
    \item Построена формальная модель конвейера «захват кадра -- детекция -- сопровождение -- правила -- исходящие сообщения», включающая декомпозицию сквозной задержки, дедупликацию событий и критерии архитектурной приемки.
    \item Реализована программная система в виде набора сервисов \texttt{api}, \texttt{analytics-worker}, \texttt{integration-worker}, \texttt{ui} с инфраструктурой PostgreSQL, RabbitMQ, Prometheus, Grafana и Docker Compose.
    \item Проведена стендовая апробация: подтверждены интеграционная доставка синтетически зафиксированного события \texttt{zone\_enter}, запрет операции записи для роли \texttt{viewer}, восстановление публикации после отказа RabbitMQ, повторная доставка после недоступности HTTP-получателя и связность демонстрационного БПЛА-видеосценария.
    \item Определены ограничения текущей реализации и направления развития: длительная многокамерная оценка задержки, расширение набора классов, межкамерная аналитика, объектовое хранилище артефактов и промышленная отказоустойчивая конфигурация.
\end{enumerate}

Практическая значимость работы заключается в том, что полученная платформа может использоваться как воспроизводимый исследовательско-экспериментальный стенд для задач охранного видеонаблюдения. Она позволяет проверять не только качество детектора, но и системные свойства, важные для эксплуатации: сохранность событий при отказах, управляемость повторной доставки, наблюдаемость и соблюдение ролевой модели.

Соответствие цели и задачам исследования подтверждается следующим:
\begin{enumerate}
    \item первая задача выполнена: требования к платформе сформированы на основе анализа предметной области, существующих решений и ограничений охранного видеонаблюдения;
    \item вторая задача выполнена: методы обнаружения, сопровождения и событийного анализа обоснованы сравнительным анализом и реализованы в программном контуре;
    \item третья задача выполнена: модель, архитектура и алгоритмический конвейер платформы формализованы и связаны с требованиями к задержке, доставке событий и наблюдаемости;
    \item четвертая задача выполнена: реализованы компоненты управления, аналитической обработки, интеграционной доставки, разграничения доступа, мониторинга и воспроизводимого развертывания;
    \item пятая задача выполнена: стендовая апробация зафиксирована в виде протокола RUN-2026-05-13-A1, JSON-артефакта, чек-листов, акта стендового прогона, скриншотов, графиков обучения и таблиц метрик; сценарий перехода в очередь недоставленных сообщений подтвержден отдельной проверкой логики фонового обработчика; количественные выводы ограничены фактически выполненным приемочным контуром и не подменяют промышленную p95-оценку полного видеоконтура;
    \item шестая задача выполнена: ограничения текущей реализации и направления перехода к промышленному применению явно сформулированы.
\end{enumerate}

Апробация подтверждает выполнение цели именно на уровне исследовательско-экспериментального стенда: проверены фиксация и доставка событий, восстановление после кратковременных отказов интеграционного контура, соблюдение ролевых ограничений и связность демонстрационного видеосценария. Эти результаты достаточны для вывода о работоспособности архитектурно-алгоритмического подхода, но не являются заявкой на промышленную сертификацию платформы.

Ограничения работы также зафиксированы явно: текущая количественная апробация выполнена на CPU-профиле приемки для интеграционного контура, демонстрационный набор БПЛА подтверждает связность видеосценария, а промышленная многокамерная p95-оценка полного пути от объекта к событию требует отдельной длительной серии прогонов с экспортом временных рядов Prometheus/Grafana. Сценарий длительной недоступности HTTP-получателя до исчерпания повторов не измерялся в RUN-2026-05-13-A1 и отнесен к следующему стендовому этапу. Кроме того, текущая версия не включает межкамерную реидентификацию, кластерную отказоустойчивость PostgreSQL/RabbitMQ и промышленную калибровку детектора на объектовых сценах.

Дальнейшее развитие целесообразно вести по следующим направлениям:
\begin{itemize}
    \item расширение набора поддерживаемых классов и дообучение детектора на предметных наборах данных конкретных охраняемых объектов;
    \item проведение длительной многокамерной серии прогонов для статистически устойчивой оценки \(\operatorname{p95}_{\text{obj}\to\text{event}}\);
    \item внедрение межкамерной корреляции и реидентификации;
    \item перенос снимков и видеоартефактов в объектное хранилище;
    \item кластеризация брокера и базы данных для эксплуатационного профиля высокой доступности;
    \item аппаратное ускорение инференса через ONNX Runtime или TensorRT для плотных многокамерных конфигураций.
\end{itemize}

Таким образом, поставленная цель достигнута на уровне исследовательско-экспериментального стенда: разработана, реализована и апробирована платформа видеоаналитики реального времени для мониторинга охраняемых зон и событийного анализа, обеспечивающая формирование интерпретируемых событий, их сохранение и надежную доставку сведений о событиях внешним потребителям.
